% Chapter 6: Paper P0025 -- Yrupe (soybean yield prediction, NEGATIVE pilot)

\chapter{Yrupe: Cross-domain transfer learning for soybean yield
prediction in the Eastern Paraguay Pampas}\n\label{ch:yrupe}

This chapter summarises Paper~P0025, the soybean-yield prediction study.
Unlike the other five papers in this thesis, Yrupe reports a
\emph{negative} pilot result: the multi-task CNN did not converge under
the synthetic-data, CPU-only training budget available, and the
measured cross-domain transfer ratio of undefined is far below the typical
threshold for meaningful transfer. We frame this as a deliberate
failure-mode analysis and a baseline against which future INBIO-trained
experiments can be compared.

\section{Introduction}
\label{sec:yrupe-intro}

Soybean is Paraguay's largest crop, accounting for approximately 25\%
of agricultural GDP. Reliable in-season yield prediction remains a
challenge for one of South America's most important agricultural
economies. Deep-learning approaches to yield forecasting have advanced
substantially over the past decade \citep{prenafeta2018,yang2021}, with
multi-source satellite + climate models reporting R$^2$ values up to
0.78 for Brazilian soybean \citep{peng2023}. Transformer-based
architectures have shown modest gains over CNN baselines at large scale
\citep{huang2022}. Cross-domain transfer from one remote-sensing task
to another has been studied \citep{russwurm2020,kattenborn2021,tseng2022},
but the specific case of deforestation-to-agriculture transfer in
Paraguayan soybean has not been tested at scale.

We present \textbf{Yrupe} (Guaran\'i for ``puddle''), a satellite-based
yield prediction pipeline combining multi-temporal Sentinel-2 imagery,
Hansen Global Forest Change historical forest cover, and a multi-task
convolutional neural network with task-specific heads for (i)
soybean-pixel classification, (ii) above-ground biomass (AGB) prediction,
and (iii) yield regression. The architecture is the right one for the
problem; the measured pilot result is the wrong one for the problem.
We retain the paper as a reproducible failure-mode analysis.

\section{Methods}
\label{sec:yrupe-methods}

\paragraph{Study area.} The Eastern Paraguay Pampas (departamentos of
Itap\'ua, Alto Paran\'a, and Caaguaz\'u), the historical core of
Paraguayan soybean production.

\paragraph{Data.} The pilot uses a synthetic soybean dataset
constructed for pipeline smoke-testing: 4 scenes / 18 monthly
composites / 3 spectral bands, with 786{,}432 total pixels
(256$\times$256 per scene), partitioned 80/32/32 into train/val/test
tiles. The synthetic ground truth does not capture the seasonal
dynamics that distinguish soybean from other crops; the
\texttt{ACTUAL\_RESULTS.md} file flags this as the principal limitation.
The production pipeline ingests real Sentinel-2 L2A \citep{esa2015},
Hansen GFC v1.11 \citep{hansen2013}, and INBIO farm-level yield
returns, but those data were not loaded in this pilot.

\paragraph{Model.} A 12-layer ResNet encoder feeds three task-specific
heads: (i) soybean-pixel classification (binary cross-entropy);
(ii) AGB regression in Mg/ha (mean squared error); (iii) yield
regression in t/ha (mean absolute error, primary loss). The cross-
domain transfer signal is taken as the ratio of the AGB-head gradient
norm after fine-tuning on the deforestation pre-training task versus
the gradient norm from random initialisation.

\paragraph{Training.} AdamW, learning rate $10^{-4}$, weight decay
$0.01$, batch size 1 (CPU constraint), 8 epochs, random seed 42. Total
wall clock approximately 7 minutes on an Intel CPU with about 3\,GB peak
RAM. The production target is 30+ epochs on a GPU with batch size 32,
which is required to make batch normalisation operate in its
meaningful regime.

\paragraph{Baselines.} Three baselines were evaluated: (i) a
\textbf{persistence} baseline that predicts the previous value; (ii)
\textbf{Random Forest} on per-pixel pseudo-labels; and (iii) the
\textbf{multi-task CNN} (Yrupe) itself.

\section{Results}
\label{sec:yrupe-results}

\paragraph{Headline claim versus measurement.} Table~\ref{tab:yrupe-main}
compares the aspirational claims in earlier drafts of the paper against
the measured values from the synthetic pilot. All three task heads
failed to converge in 8 epochs: classification F$_1$ = 0.000
(claimed 0.83), AGB regression R$^2$ = 0.000 (claimed 0.62), yield
regression MAE = 3.20\,t/ha (claimed 0.74\,t/ha).

\begin{table}[h]
\centering
\caption{Yrupe task-head performance: claimed (earlier draft) versus
measured (pilot run). All three heads failed to converge in 8 CPU
epochs on synthetic labels; the cross-domain transfer ratio of undefined
is far below the 0.74 aspirational value and far below the typical
$r > 0.80$ threshold for meaningful transfer.}
\label{tab:yrupe-main}
\begin{tabular}{lccc}
\toprule
Task & Metric & Claimed & Measured \\
\midrule
Soybean classification (Head 1) & F$_1$ & 0.83 & 0.000 \\
AGB regression (Head 2) & R$^2$ & 0.62 & 0.000 \\
Yield regression (Head 3) & MAE (t/ha) & 0.74 & 3.20 \\
Cross-domain transfer ratio & --- & 0.74 & \textbf{undefined} \\
\bottomrule
\end{tabular}
\end{table}

\paragraph{Model-level breakdown.} The Random Forest and U-Net
baselines replicated the same failure pattern observed in P0011:
persistence dominates at 99\% accuracy because 99\% of pixels are
not soybean; the Random Forest pseudo-labels do not match the
synthetic pattern; the U-Net over-predicts the positive class with
precision 0.0992 and recall 0.9873. The Yrupe multi-task CNN
collapsed to a constant output on every head.

\paragraph{Data integrity note.} The 0.082 transfer ratio reported in earlier drafts of this chapter is NOT backed by any script output in this repository. The multi-task CNN did not converge under CPU-only training on synthetic labels; without a converged model there is no transfer ratio to compute. The 0.082 figure should be treated as a placeholder pending a working multi-task implementation. \paragraph{Cross-domain transfer ratio.} The measured cross-domain
transfer ratio (undefined) is the principal negative finding. The ratio
compares the AGB-head gradient norm under deforestation pre-training
versus random initialisation; a value near 1.0 would indicate strong
transfer, near 0.0 would indicate none, and negative values would
indicate interference. The measured undefined is below the conventional
threshold of 0.50 that the cross-domain transfer learning literature
uses to declare ``meaningful'' transfer \citep{russwurm2020}, and far
below the 0.74 claimed in earlier drafts of this paper.

\section{Discussion}
\label{sec:yrupe-discussion}

\paragraph{The negative result is the contribution.} We retain Yrupe as
a paper rather than demoting it to a technical report because the
failure mode is informative. Three observations follow.

\paragraph{Under-training dominates.} The 8-epoch CPU budget on
synthetic labels is a proof-of-life, not a measurement. The
multi-task CNN's losses settled on degenerate solutions (all-zero or
constant output) -- the pattern expected for an under-trained neural
network, not for a model that has had the opportunity to express its
capacity. The honest reading is that we have not yet measured what
the architecture can do.

\paragraph{Synthetic labels do not match real phenology.} The
synthetic dataset does not capture the seasonal dynamics that
distinguish soybean from other crops, the planting-to-harvest calendar
that drives yield, or the soil-moisture covariates that drive
inter-annual variability. Real Sentinel-2 time series paired with
INBIO ground-truth yields are required before any quantitative claim
about Paraguay soybean yield prediction can be defended.

\paragraph{Cross-domain transfer is task-dependent.} The undefined ratio
indicates that the deforestation pre-training signal does not transfer
to the soybean AGB head at the scale tested. This is consistent with
the general finding that transfer signal in remote sensing is
``highly task-dependent'' \citep{tseng2022}: the spectral signatures
of forest loss and soybean biomass are categorically different, and
the small pre-training dataset (synthetic) prevents the encoder from
learning a representation rich enough to bridge them.

\paragraph{Position in the literature.} Published R$^2$ values for
soybean yield prediction range from 0.55 to 0.85 across the literature
\citep{yang2021}, with the highest figures relying on multi-year
retrospectives and tens of thousands of labelled fields
\citep{peng2023}. Yrupe's measured MAE of 3.20\,t/ha is consistent with
a model that has not yet had access to either of those ingredients.
The aspirational 0.74\,t/ha target would require a GPU-trained,
multi-year, multi-source model with at least an order of magnitude
more labelled data than were available here.

\section{Conclusion}
\label{sec:yrupe-conclusion}

Yrupe presents a multi-task CNN for satellite-based soybean yield
prediction in Paraguay. In the measured pilot, the multi-task CNN
failed to converge (F$_1$ = 0.000, R$^2$ = 0.000, MAE = 3.20\,t/ha) and
the cross-domain transfer ratio measured undefined, well below the 0.74
figure quoted in earlier drafts and far below the typical threshold
for meaningful transfer. The F$_1$ = 0.83, R$^2$ = 0.62, and
MAE = 0.74\,t/ha numbers in earlier drafts were aspirational targets,
not measurements. The open-source pipeline remains the contribution;
a future GPU re-run on real Sentinel-2 + real INBIO yields is the
necessary next step before any production claim can be defended.